\section{清初证据链事件锚点}
以下锚点补入清初账本中关于前史、文件、决策与后续的独立记录，以显示征服、海疆、驻防和地方治理并非由一份军报或诏令即可完成。
\subsection{1640—1649}
\eventanchor{EQ-1644-SHANHAIGUAN-PRELUDE}{「清军与吴三桂在山海关击败李自…}顺治元年四月（1644年5月，换算待核），「清军与吴三桂在山海关击败李自成军，打开进入华北的军事通道」之前的条件、信息和争议已通过中央与地方材料显现，构成该事件可追踪的前史节点。核心取证：《清世祖实录》顺治元年四月山海关军报；《清史稿·世祖本纪》；《清初内国史院满文档案》顺治元年四月关宁军务档；同时期地方档案、地方志、日记或对方材料应按具体地区补检。这一锚点界定可回查的前史条件；条件、传闻或信息累积不应被写成已经证实的单一直接原因。来源保留：官修编年和地方材料的记录密度与立场并不相同；本行不将条件、传闻、呈报或后见解释直接等同为确定因果。
\eventanchor{EQ-1644-SHANHAIGUAN-DECISION}{围绕「清军与吴三桂在山海关击败…}顺治元年四月（1644年5月，换算待核），围绕「清军与吴三桂在山海关击败李自成军，打开进入华北的军事通道」的命令、调度、照会或呈报形成独立的中央—地方行动文书链。核心取证：《清世祖实录》顺治元年四月山海关军报；《清史稿·世祖本纪》；《清初内国史院满文档案》顺治元年四月关宁军务档。这一锚点定位决策或调度动作；命令发出、地方接收、执行过程和实际结果仍须分别寻找证据。来源保留：奏折、上谕、部议、军机档或条约可证明行政动作和机构立场，不自动证明所有地区、群体或对手按其意图行动。
\eventanchor{EQ-1644-SHANHAIGUAN-AFTERMATH}{「清军与吴三桂在山海关击败李自…}顺治元年四月（1644年5月，换算待核），「清军与吴三桂在山海关击败李自成军，打开进入华北的军事通道」引发的善后、执行或地方回应构成可由不同材料追踪的后续节点。核心取证：《清世祖实录》顺治元年四月山海关军报；《清史稿·世祖本纪》；《清初内国史院满文档案》顺治元年四月关宁军务档；相关地区的档案、志书、契约、报刊或对方材料。这一锚点追踪后续影响；不同地区、群体与时段的反应不能由战后或改革后的概括性记忆代替。来源保留：战后、改革后或条约后的记忆、地方记录和官方总结常具有事后选择性；后续影响须按地点、群体和时间分层，不作整体化推断。
\eventanchor{EQ-1644-BEIJING-ENTRY-PRELUDE}{「多尔衮率军进入北京并接收明朝…}顺治元年五月（1644年6月前后），「多尔衮率军进入北京并接收明朝京城的行政与礼制空间」之前的条件、信息和争议已通过中央与地方材料显现，构成该事件可追踪的前史节点。核心取证：《清世祖实录》顺治元年五月；《清史稿·世祖本纪》；《甲申传信录》及北京地方记事；同时期地方档案、地方志、日记或对方材料应按具体地区补检。这一锚点界定可回查的前史条件；条件、传闻或信息累积不应被写成已经证实的单一直接原因。来源保留：官修编年和地方材料的记录密度与立场并不相同；本行不将条件、传闻、呈报或后见解释直接等同为确定因果。
\eventanchor{EQ-1644-BEIJING-ENTRY-DECISION}{围绕「多尔衮率军进入北京并接收…}顺治元年五月（1644年6月前后），围绕「多尔衮率军进入北京并接收明朝京城的行政与礼制空间」的命令、调度、照会或呈报形成独立的中央—地方行动文书链。核心取证：《清世祖实录》顺治元年五月；《清史稿·世祖本纪》；《甲申传信录》及北京地方记事。这一锚点定位决策或调度动作；命令发出、地方接收、执行过程和实际结果仍须分别寻找证据。来源保留：奏折、上谕、部议、军机档或条约可证明行政动作和机构立场，不自动证明所有地区、群体或对手按其意图行动。
\eventanchor{EQ-1644-BEIJING-ENTRY-AFTERMATH}{「多尔衮率军进入北京并接收明朝…}顺治元年五月（1644年6月前后），「多尔衮率军进入北京并接收明朝京城的行政与礼制空间」引发的善后、执行或地方回应构成可由不同材料追踪的后续节点。核心取证：《清世祖实录》顺治元年五月；《清史稿·世祖本纪》；《甲申传信录》及北京地方记事；相关地区的档案、志书、契约、报刊或对方材料。这一锚点追踪后续影响；不同地区、群体与时段的反应不能由战后或改革后的概括性记忆代替。来源保留：战后、改革后或条约后的记忆、地方记录和官方总结常具有事后选择性；后续影响须按地点、群体和时间分层，不作整体化推断。
\eventanchor{EQ-1644-SHUNZHI-BEIJING-PRELUDE}{「顺治帝自盛京抵北京，清廷举行…}顺治元年九月（1644年10月前后），「顺治帝自盛京抵北京，清廷举行定都与皇帝入城仪式」之前的条件、信息和争议已通过中央与地方材料显现，构成该事件可追踪的前史节点。核心取证：《清世祖实录》顺治元年九月；《清史稿·世祖本纪》；《清初内国史院满文档案》北京迁都礼仪档；同时期地方档案、地方志、日记或对方材料应按具体地区补检。这一锚点界定可回查的前史条件；条件、传闻或信息累积不应被写成已经证实的单一直接原因。来源保留：官修编年和地方材料的记录密度与立场并不相同；本行不将条件、传闻、呈报或后见解释直接等同为确定因果。
\eventanchor{EQ-1644-SHUNZHI-BEIJING-DECISION}{围绕「顺治帝自盛京抵北京，清廷…}顺治元年九月（1644年10月前后），围绕「顺治帝自盛京抵北京，清廷举行定都与皇帝入城仪式」的命令、调度、照会或呈报形成独立的中央—地方行动文书链。核心取证：《清世祖实录》顺治元年九月；《清史稿·世祖本纪》；《清初内国史院满文档案》北京迁都礼仪档。这一锚点定位决策或调度动作；命令发出、地方接收、执行过程和实际结果仍须分别寻找证据。来源保留：奏折、上谕、部议、军机档或条约可证明行政动作和机构立场，不自动证明所有地区、群体或对手按其意图行动。
\eventanchor{EQ-1644-SHUNZHI-BEIJING-AFTERMATH}{「顺治帝自盛京抵北京，清廷举行…}顺治元年九月（1644年10月前后），「顺治帝自盛京抵北京，清廷举行定都与皇帝入城仪式」引发的善后、执行或地方回应构成可由不同材料追踪的后续节点。核心取证：《清世祖实录》顺治元年九月；《清史稿·世祖本纪》；《清初内国史院满文档案》北京迁都礼仪档；相关地区的档案、志书、契约、报刊或对方材料。这一锚点追踪后续影响；不同地区、群体与时段的反应不能由战后或改革后的概括性记忆代替。来源保留：战后、改革后或条约后的记忆、地方记录和官方总结常具有事后选择性；后续影响须按地点、群体和时间分层，不作整体化推断。
\eventanchor{EQ-1644-REGENT-ORDERS-PRELUDE}{「多尔衮以摄政王名义整顿京师官…}顺治元年冬（1644年），「多尔衮以摄政王名义整顿京师官署、招降明官并发布安民与征发命令」之前的条件、信息和争议已通过中央与地方材料显现，构成该事件可追踪的前史节点。核心取证：《清世祖实录》顺治元年十月至十二月；《清史稿·职官志》；中国第一历史档案馆藏内阁大库顺治元年题本〔京师官署交接〕；同时期地方档案、地方志、日记或对方材料应按具体地区补检。这一锚点界定可回查的前史条件；条件、传闻或信息累积不应被写成已经证实的单一直接原因。来源保留：官修编年和地方材料的记录密度与立场并不相同；本行不将条件、传闻、呈报或后见解释直接等同为确定因果。
\eventanchor{EQ-1644-REGENT-ORDERS-DECISION}{围绕「多尔衮以摄政王名义整顿京…}顺治元年冬（1644年），围绕「多尔衮以摄政王名义整顿京师官署、招降明官并发布安民与征发命令」的命令、调度、照会或呈报形成独立的中央—地方行动文书链。核心取证：《清世祖实录》顺治元年十月至十二月；《清史稿·职官志》；中国第一历史档案馆藏内阁大库顺治元年题本〔京师官署交接〕。这一锚点定位决策或调度动作；命令发出、地方接收、执行过程和实际结果仍须分别寻找证据。来源保留：奏折、上谕、部议、军机档或条约可证明行政动作和机构立场，不自动证明所有地区、群体或对手按其意图行动。
\eventanchor{EQ-1644-REGENT-ORDERS-AFTERMATH}{「多尔衮以摄政王名义整顿京师官…}顺治元年冬（1644年），「多尔衮以摄政王名义整顿京师官署、招降明官并发布安民与征发命令」引发的善后、执行或地方回应构成可由不同材料追踪的后续节点。核心取证：《清世祖实录》顺治元年十月至十二月；《清史稿·职官志》；中国第一历史档案馆藏内阁大库顺治元年题本〔京师官署交接〕；相关地区的档案、志书、契约、报刊或对方材料。这一锚点追踪后续影响；不同地区、群体与时段的反应不能由战后或改革后的概括性记忆代替。来源保留：战后、改革后或条约后的记忆、地方记录和官方总结常具有事后选择性；后续影响须按地点、群体和时间分层，不作整体化推断。
\eventanchor{EQ-1645-NANJING-TAKEN-PRELUDE}{「多铎军攻取南京，南明弘光朝的…}顺治二年四月（1645年5月前后），「多铎军攻取南京，南明弘光朝的都城失守」之前的条件、信息和争议已通过中央与地方材料显现，构成该事件可追踪的前史节点。核心取证：《清世祖实录》顺治二年四月；《清史稿·世祖本纪》；《南明史料》所收弘光朝奏疏与江宁地方记事；同时期地方档案、地方志、日记或对方材料应按具体地区补检。这一锚点界定可回查的前史条件；条件、传闻或信息累积不应被写成已经证实的单一直接原因。来源保留：官修编年和地方材料的记录密度与立场并不相同；本行不将条件、传闻、呈报或后见解释直接等同为确定因果。
\eventanchor{EQ-1645-NANJING-TAKEN-DECISION}{围绕「多铎军攻取南京，南明弘光…}顺治二年四月（1645年5月前后），围绕「多铎军攻取南京，南明弘光朝的都城失守」的命令、调度、照会或呈报形成独立的中央—地方行动文书链。核心取证：《清世祖实录》顺治二年四月；《清史稿·世祖本纪》；《南明史料》所收弘光朝奏疏与江宁地方记事。这一锚点定位决策或调度动作；命令发出、地方接收、执行过程和实际结果仍须分别寻找证据。来源保留：奏折、上谕、部议、军机档或条约可证明行政动作和机构立场，不自动证明所有地区、群体或对手按其意图行动。
\eventanchor{EQ-1645-NANJING-TAKEN-AFTERMATH}{「多铎军攻取南京，南明弘光朝的…}顺治二年四月（1645年5月前后），「多铎军攻取南京，南明弘光朝的都城失守」引发的善后、执行或地方回应构成可由不同材料追踪的后续节点。核心取证：《清世祖实录》顺治二年四月；《清史稿·世祖本纪》；《南明史料》所收弘光朝奏疏与江宁地方记事；相关地区的档案、志书、契约、报刊或对方材料。这一锚点追踪后续影响；不同地区、群体与时段的反应不能由战后或改革后的概括性记忆代替。来源保留：战后、改革后或条约后的记忆、地方记录和官方总结常具有事后选择性；后续影响须按地点、群体和时间分层，不作整体化推断。
\eventanchor{EQ-1645-HONGGUANG-CAPTURE-PRELUDE}{「弘光帝朱由崧在芜湖一带被俘并…}顺治二年五月（1645年），「弘光帝朱由崧在芜湖一带被俘并被押往北京」之前的条件、信息和争议已通过中央与地方材料显现，构成该事件可追踪的前史节点。核心取证：《清世祖实录》顺治二年五月；《清史稿·弘光帝本纪》；《明季南略》及江南地方志相关条目；同时期地方档案、地方志、日记或对方材料应按具体地区补检。这一锚点界定可回查的前史条件；条件、传闻或信息累积不应被写成已经证实的单一直接原因。来源保留：官修编年和地方材料的记录密度与立场并不相同；本行不将条件、传闻、呈报或后见解释直接等同为确定因果。
\eventanchor{EQ-1645-HONGGUANG-CAPTURE-DECISION}{围绕「弘光帝朱由崧在芜湖一带被…}顺治二年五月（1645年），围绕「弘光帝朱由崧在芜湖一带被俘并被押往北京」的命令、调度、照会或呈报形成独立的中央—地方行动文书链。核心取证：《清世祖实录》顺治二年五月；《清史稿·弘光帝本纪》；《明季南略》及江南地方志相关条目。这一锚点定位决策或调度动作；命令发出、地方接收、执行过程和实际结果仍须分别寻找证据。来源保留：奏折、上谕、部议、军机档或条约可证明行政动作和机构立场，不自动证明所有地区、群体或对手按其意图行动。
\eventanchor{EQ-1645-HONGGUANG-CAPTURE-AFTERMATH}{「弘光帝朱由崧在芜湖一带被俘并…}顺治二年五月（1645年），「弘光帝朱由崧在芜湖一带被俘并被押往北京」引发的善后、执行或地方回应构成可由不同材料追踪的后续节点。核心取证：《清世祖实录》顺治二年五月；《清史稿·弘光帝本纪》；《明季南略》及江南地方志相关条目；相关地区的档案、志书、契约、报刊或对方材料。这一锚点追踪后续影响；不同地区、群体与时段的反应不能由战后或改革后的概括性记忆代替。来源保留：战后、改革后或条约后的记忆、地方记录和官方总结常具有事后选择性；后续影响须按地点、群体和时间分层，不作整体化推断。
\eventanchor{EQ-1645-HAIR-ORDER-PRELUDE}{「清廷重申薙发易服命令，在江南…}顺治二年六月（1645年），「清廷重申薙发易服命令，在江南等地引发服从、逃亡与抵抗」之前的条件、信息和争议已通过中央与地方材料显现，构成该事件可追踪的前史节点。核心取证：《清世祖实录》顺治二年六月；《清史稿·舆服志》；内阁大库顺治二年兵科题本与江南府县志薙发条；同时期地方档案、地方志、日记或对方材料应按具体地区补检。这一锚点界定可回查的前史条件；条件、传闻或信息累积不应被写成已经证实的单一直接原因。来源保留：官修编年和地方材料的记录密度与立场并不相同；本行不将条件、传闻、呈报或后见解释直接等同为确定因果。
\eventanchor{EQ-1645-HAIR-ORDER-DECISION}{围绕「清廷重申薙发易服命令，在…}顺治二年六月（1645年），围绕「清廷重申薙发易服命令，在江南等地引发服从、逃亡与抵抗」的命令、调度、照会或呈报形成独立的中央—地方行动文书链。核心取证：《清世祖实录》顺治二年六月；《清史稿·舆服志》；内阁大库顺治二年兵科题本与江南府县志薙发条。这一锚点定位决策或调度动作；命令发出、地方接收、执行过程和实际结果仍须分别寻找证据。来源保留：奏折、上谕、部议、军机档或条约可证明行政动作和机构立场，不自动证明所有地区、群体或对手按其意图行动。
\eventanchor{EQ-1645-HAIR-ORDER-AFTERMATH}{「清廷重申薙发易服命令，在江南…}顺治二年六月（1645年），「清廷重申薙发易服命令，在江南等地引发服从、逃亡与抵抗」引发的善后、执行或地方回应构成可由不同材料追踪的后续节点。核心取证：《清世祖实录》顺治二年六月；《清史稿·舆服志》；内阁大库顺治二年兵科题本与江南府县志薙发条；相关地区的档案、志书、契约、报刊或对方材料。这一锚点追踪后续影响；不同地区、群体与时段的反应不能由战后或改革后的概括性记忆代替。来源保留：战后、改革后或条约后的记忆、地方记录和官方总结常具有事后选择性；后续影响须按地点、群体和时间分层，不作整体化推断。
\eventanchor{EQ-1645-YANGZHOU-PRELUDE}{「扬州攻陷后发生大规模杀戮与掳…}顺治二年四—五月（1645年），「扬州攻陷后发生大规模杀戮与掳掠，规模和持续时间存在重大争议」之前的条件、信息和争议已通过中央与地方材料显现，构成该事件可追踪的前史节点。核心取证：《清世祖实录》顺治二年四月扬州军报；《清史稿·史可法传》；王秀楚《扬州十日记》与《扬州府志》；同时期地方档案、地方志、日记或对方材料应按具体地区补检。这一锚点界定可回查的前史条件；条件、传闻或信息累积不应被写成已经证实的单一直接原因。来源保留：官修编年和地方材料的记录密度与立场并不相同；本行不将条件、传闻、呈报或后见解释直接等同为确定因果。
\eventanchor{EQ-1645-YANGZHOU-DECISION}{围绕「扬州攻陷后发生大规模杀戮…}顺治二年四—五月（1645年），围绕「扬州攻陷后发生大规模杀戮与掳掠，规模和持续时间存在重大争议」的命令、调度、照会或呈报形成独立的中央—地方行动文书链。核心取证：《清世祖实录》顺治二年四月扬州军报；《清史稿·史可法传》；王秀楚《扬州十日记》与《扬州府志》。这一锚点定位决策或调度动作；命令发出、地方接收、执行过程和实际结果仍须分别寻找证据。来源保留：奏折、上谕、部议、军机档或条约可证明行政动作和机构立场，不自动证明所有地区、群体或对手按其意图行动。
\eventanchor{EQ-1645-YANGZHOU-AFTERMATH}{「扬州攻陷后发生大规模杀戮与掳…}顺治二年四—五月（1645年），「扬州攻陷后发生大规模杀戮与掳掠，规模和持续时间存在重大争议」引发的善后、执行或地方回应构成可由不同材料追踪的后续节点。核心取证：《清世祖实录》顺治二年四月扬州军报；《清史稿·史可法传》；王秀楚《扬州十日记》与《扬州府志》；相关地区的档案、志书、契约、报刊或对方材料。这一锚点追踪后续影响；不同地区、群体与时段的反应不能由战后或改革后的概括性记忆代替。来源保留：战后、改革后或条约后的记忆、地方记录和官方总结常具有事后选择性；后续影响须按地点、群体和时间分层，不作整体化推断。
\eventanchor{EQ-1645-JIANGYIN-PRELUDE}{「江阴在薙发命令和清军围攻中失…}顺治二年闰六月至八月（1645年），「江阴在薙发命令和清军围攻中失守，地方记忆长期以守城与屠城叙述该事件」之前的条件、信息和争议已通过中央与地方材料显现，构成该事件可追踪的前史节点。核心取证：《清世祖实录》顺治二年江南军报；《清史稿·世祖本纪》；《江阴县志》及《江阴城守纪》；同时期地方档案、地方志、日记或对方材料应按具体地区补检。这一锚点界定可回查的前史条件；条件、传闻或信息累积不应被写成已经证实的单一直接原因。来源保留：官修编年和地方材料的记录密度与立场并不相同；本行不将条件、传闻、呈报或后见解释直接等同为确定因果。
\eventanchor{EQ-1645-JIANGYIN-DECISION}{围绕「江阴在薙发命令和清军围攻…}顺治二年闰六月至八月（1645年），围绕「江阴在薙发命令和清军围攻中失守，地方记忆长期以守城与屠城叙述该事件」的命令、调度、照会或呈报形成独立的中央—地方行动文书链。核心取证：《清世祖实录》顺治二年江南军报；《清史稿·世祖本纪》；《江阴县志》及《江阴城守纪》。这一锚点定位决策或调度动作；命令发出、地方接收、执行过程和实际结果仍须分别寻找证据。来源保留：奏折、上谕、部议、军机档或条约可证明行政动作和机构立场，不自动证明所有地区、群体或对手按其意图行动。
\eventanchor{EQ-1645-JIANGYIN-AFTERMATH}{「江阴在薙发命令和清军围攻中失…}顺治二年闰六月至八月（1645年），「江阴在薙发命令和清军围攻中失守，地方记忆长期以守城与屠城叙述该事件」引发的善后、执行或地方回应构成可由不同材料追踪的后续节点。核心取证：《清世祖实录》顺治二年江南军报；《清史稿·世祖本纪》；《江阴县志》及《江阴城守纪》；相关地区的档案、志书、契约、报刊或对方材料。这一锚点追踪后续影响；不同地区、群体与时段的反应不能由战后或改革后的概括性记忆代替。来源保留：战后、改革后或条约后的记忆、地方记录和官方总结常具有事后选择性；后续影响须按地点、群体和时间分层，不作整体化推断。
\eventanchor{EQ-1645-ZHEJIANG-CONQUEST-PRELUDE}{「清军进取杭州、绍兴等地，南明…}顺治二年六月至十月（1645年），「清军进取杭州、绍兴等地，南明鲁王政权转入浙东及海上活动」之前的条件、信息和争议已通过中央与地方材料显现，构成该事件可追踪的前史节点。核心取证：《清世祖实录》顺治二年六月至十月；《清史稿·鲁王传》；《浙江通志》与鲁王政权文书辑录；同时期地方档案、地方志、日记或对方材料应按具体地区补检。这一锚点界定可回查的前史条件；条件、传闻或信息累积不应被写成已经证实的单一直接原因。来源保留：官修编年和地方材料的记录密度与立场并不相同；本行不将条件、传闻、呈报或后见解释直接等同为确定因果。
\eventanchor{EQ-1645-ZHEJIANG-CONQUEST-DECISION}{围绕「清军进取杭州、绍兴等地，…}顺治二年六月至十月（1645年），围绕「清军进取杭州、绍兴等地，南明鲁王政权转入浙东及海上活动」的命令、调度、照会或呈报形成独立的中央—地方行动文书链。核心取证：《清世祖实录》顺治二年六月至十月；《清史稿·鲁王传》；《浙江通志》与鲁王政权文书辑录。这一锚点定位决策或调度动作；命令发出、地方接收、执行过程和实际结果仍须分别寻找证据。来源保留：奏折、上谕、部议、军机档或条约可证明行政动作和机构立场，不自动证明所有地区、群体或对手按其意图行动。
\eventanchor{EQ-1645-ZHEJIANG-CONQUEST-AFTERMATH}{「清军进取杭州、绍兴等地，南明…}顺治二年六月至十月（1645年），「清军进取杭州、绍兴等地，南明鲁王政权转入浙东及海上活动」引发的善后、执行或地方回应构成可由不同材料追踪的后续节点。核心取证：《清世祖实录》顺治二年六月至十月；《清史稿·鲁王传》；《浙江通志》与鲁王政权文书辑录；相关地区的档案、志书、契约、报刊或对方材料。这一锚点追踪后续影响；不同地区、群体与时段的反应不能由战后或改革后的概括性记忆代替。来源保留：战后、改革后或条约后的记忆、地方记录和官方总结常具有事后选择性；后续影响须按地点、群体和时间分层，不作整体化推断。
\eventanchor{EQ-1646-LONGWU-CAPTURE-PRELUDE}{「清军在闽北俘杀隆武帝朱聿键，…}顺治三年八月（1646年），「清军在闽北俘杀隆武帝朱聿键，隆武朝终结」之前的条件、信息和争议已通过中央与地方材料显现，构成该事件可追踪的前史节点。核心取证：《清世祖实录》顺治三年八月；《清史稿·隆武帝本纪》；《福建通志》与黄道周文集相关记载；同时期地方档案、地方志、日记或对方材料应按具体地区补检。这一锚点界定可回查的前史条件；条件、传闻或信息累积不应被写成已经证实的单一直接原因。来源保留：官修编年和地方材料的记录密度与立场并不相同；本行不将条件、传闻、呈报或后见解释直接等同为确定因果。
\eventanchor{EQ-1646-LONGWU-CAPTURE-DECISION}{围绕「清军在闽北俘杀隆武帝朱聿…}顺治三年八月（1646年），围绕「清军在闽北俘杀隆武帝朱聿键，隆武朝终结」的命令、调度、照会或呈报形成独立的中央—地方行动文书链。核心取证：《清世祖实录》顺治三年八月；《清史稿·隆武帝本纪》；《福建通志》与黄道周文集相关记载。这一锚点定位决策或调度动作；命令发出、地方接收、执行过程和实际结果仍须分别寻找证据。来源保留：奏折、上谕、部议、军机档或条约可证明行政动作和机构立场，不自动证明所有地区、群体或对手按其意图行动。
\eventanchor{EQ-1646-LONGWU-CAPTURE-AFTERMATH}{「清军在闽北俘杀隆武帝朱聿键，…}顺治三年八月（1646年），「清军在闽北俘杀隆武帝朱聿键，隆武朝终结」引发的善后、执行或地方回应构成可由不同材料追踪的后续节点。核心取证：《清世祖实录》顺治三年八月；《清史稿·隆武帝本纪》；《福建通志》与黄道周文集相关记载；相关地区的档案、志书、契约、报刊或对方材料。这一锚点追踪后续影响；不同地区、群体与时段的反应不能由战后或改革后的概括性记忆代替。来源保留：战后、改革后或条约后的记忆、地方记录和官方总结常具有事后选择性；后续影响须按地点、群体和时间分层，不作整体化推断。
\eventanchor{EQ-1646-SHAOWU-GUANGZHOU-PRELUDE}{「清军攻入广州，绍武帝绍武帝…}顺治三年十一月（1646年），「清军攻入广州，绍武帝绍武帝被杀，广州短暂的绍武朝结束」之前的条件、信息和争议已通过中央与地方材料显现，构成该事件可追踪的前史节点。核心取证：《清世祖实录》顺治三年十一月；《清史稿·绍武帝本纪》；《广东通志》及南明文集；同时期地方档案、地方志、日记或对方材料应按具体地区补检。这一锚点界定可回查的前史条件；条件、传闻或信息累积不应被写成已经证实的单一直接原因。来源保留：官修编年和地方材料的记录密度与立场并不相同；本行不将条件、传闻、呈报或后见解释直接等同为确定因果。
\eventanchor{EQ-1646-SHAOWU-GUANGZHOU-DECISION}{围绕「清军攻入广州，绍武帝朱聿…}顺治三年十一月（1646年），围绕「清军攻入广州，绍武帝绍武帝被杀，广州短暂的绍武朝结束」的命令、调度、照会或呈报形成独立的中央—地方行动文书链。核心取证：《清世祖实录》顺治三年十一月；《清史稿·绍武帝本纪》；《广东通志》及南明文集。这一锚点定位决策或调度动作；命令发出、地方接收、执行过程和实际结果仍须分别寻找证据。来源保留：奏折、上谕、部议、军机档或条约可证明行政动作和机构立场，不自动证明所有地区、群体或对手按其意图行动。
\eventanchor{EQ-1646-SHAOWU-GUANGZHOU-AFTERMATH}{「清军攻入广州，绍武帝绍武帝…}顺治三年十一月（1646年），「清军攻入广州，绍武帝绍武帝被杀，广州短暂的绍武朝结束」引发的善后、执行或地方回应构成可由不同材料追踪的后续节点。核心取证：《清世祖实录》顺治三年十一月；《清史稿·绍武帝本纪》；《广东通志》及南明文集；相关地区的档案、志书、契约、报刊或对方材料。这一锚点追踪后续影响；不同地区、群体与时段的反应不能由战后或改革后的概括性记忆代替。来源保留：战后、改革后或条约后的记忆、地方记录和官方总结常具有事后选择性；后续影响须按地点、群体和时间分层，不作整体化推断。
\eventanchor{EQ-1646-YONGLI-ZHAOQING-PRELUDE}{「朱由榔在肇庆即位为永历帝，建…}永历元年／顺治三年（1646年冬），「朱由榔在肇庆即位为永历帝，建立延续至西南的南明朝廷」之前的条件、信息和争议已通过中央与地方材料显现，构成该事件可追踪的前史节点。核心取证：《清世祖实录》顺治三年两广军报；《清史稿·永历帝本纪》；《永历实录》与《肇庆府志》；同时期地方档案、地方志、日记或对方材料应按具体地区补检。这一锚点界定可回查的前史条件；条件、传闻或信息累积不应被写成已经证实的单一直接原因。来源保留：官修编年和地方材料的记录密度与立场并不相同；本行不将条件、传闻、呈报或后见解释直接等同为确定因果。
\eventanchor{EQ-1646-YONGLI-ZHAOQING-DECISION}{围绕「朱由榔在肇庆即位为永历帝…}永历元年／顺治三年（1646年冬），围绕「朱由榔在肇庆即位为永历帝，建立延续至西南的南明朝廷」的命令、调度、照会或呈报形成独立的中央—地方行动文书链。核心取证：《清世祖实录》顺治三年两广军报；《清史稿·永历帝本纪》；《永历实录》与《肇庆府志》。这一锚点定位决策或调度动作；命令发出、地方接收、执行过程和实际结果仍须分别寻找证据。来源保留：奏折、上谕、部议、军机档或条约可证明行政动作和机构立场，不自动证明所有地区、群体或对手按其意图行动。
\eventanchor{EQ-1646-YONGLI-ZHAOQING-AFTERMATH}{「朱由榔在肇庆即位为永历帝，建…}永历元年／顺治三年（1646年冬），「朱由榔在肇庆即位为永历帝，建立延续至西南的南明朝廷」引发的善后、执行或地方回应构成可由不同材料追踪的后续节点。核心取证：《清世祖实录》顺治三年两广军报；《清史稿·永历帝本纪》；《永历实录》与《肇庆府志》；相关地区的档案、志书、契约、报刊或对方材料。这一锚点追踪后续影响；不同地区、群体与时段的反应不能由战后或改革后的概括性记忆代替。来源保留：战后、改革后或条约后的记忆、地方记录和官方总结常具有事后选择性；后续影响须按地点、群体和时间分层，不作整体化推断。
\eventanchor{EQ-1648-JIN-SHENGHUAN-PRELUDE}{「南昌守将金声桓反清并改奉南明…}顺治五年（1648年），「南昌守将金声桓反清并改奉南明，江西战局重新打开」之前的条件、信息和争议已通过中央与地方材料显现，构成该事件可追踪的前史节点。核心取证：《清世祖实录》顺治五年江西军报；《清史稿·金声桓传》；《南昌府志》与南明永历朝文书；同时期地方档案、地方志、日记或对方材料应按具体地区补检。这一锚点界定可回查的前史条件；条件、传闻或信息累积不应被写成已经证实的单一直接原因。来源保留：官修编年和地方材料的记录密度与立场并不相同；本行不将条件、传闻、呈报或后见解释直接等同为确定因果。
\eventanchor{EQ-1648-JIN-SHENGHUAN-DECISION}{围绕「南昌守将金声桓反清并改奉…}顺治五年（1648年），围绕「南昌守将金声桓反清并改奉南明，江西战局重新打开」的命令、调度、照会或呈报形成独立的中央—地方行动文书链。核心取证：《清世祖实录》顺治五年江西军报；《清史稿·金声桓传》；《南昌府志》与南明永历朝文书。这一锚点定位决策或调度动作；命令发出、地方接收、执行过程和实际结果仍须分别寻找证据。来源保留：奏折、上谕、部议、军机档或条约可证明行政动作和机构立场，不自动证明所有地区、群体或对手按其意图行动。
\eventanchor{EQ-1648-JIN-SHENGHUAN-AFTERMATH}{「南昌守将金声桓反清并改奉南明…}顺治五年（1648年），「南昌守将金声桓反清并改奉南明，江西战局重新打开」引发的善后、执行或地方回应构成可由不同材料追踪的后续节点。核心取证：《清世祖实录》顺治五年江西军报；《清史稿·金声桓传》；《南昌府志》与南明永历朝文书；相关地区的档案、志书、契约、报刊或对方材料。这一锚点追踪后续影响；不同地区、群体与时段的反应不能由战后或改革后的概括性记忆代替。来源保留：战后、改革后或条约后的记忆、地方记录和官方总结常具有事后选择性；后续影响须按地点、群体和时间分层，不作整体化推断。
\eventanchor{EQ-1648-LI-CHENGDONG-PRELUDE}{「广东清将李成栋反清，广州一度…}顺治五年（1648年），「广东清将李成栋反清，广州一度转入永历朝控制」之前的条件、信息和争议已通过中央与地方材料显现，构成该事件可追踪的前史节点。核心取证：《清世祖实录》顺治五年广东军报；《清史稿·李成栋传》；《广东通志》及永历朝诏敕辑录；同时期地方档案、地方志、日记或对方材料应按具体地区补检。这一锚点界定可回查的前史条件；条件、传闻或信息累积不应被写成已经证实的单一直接原因。来源保留：官修编年和地方材料的记录密度与立场并不相同；本行不将条件、传闻、呈报或后见解释直接等同为确定因果。
\eventanchor{EQ-1648-LI-CHENGDONG-DECISION}{围绕「广东清将李成栋反清，广州…}顺治五年（1648年），围绕「广东清将李成栋反清，广州一度转入永历朝控制」的命令、调度、照会或呈报形成独立的中央—地方行动文书链。核心取证：《清世祖实录》顺治五年广东军报；《清史稿·李成栋传》；《广东通志》及永历朝诏敕辑录。这一锚点定位决策或调度动作；命令发出、地方接收、执行过程和实际结果仍须分别寻找证据。来源保留：奏折、上谕、部议、军机档或条约可证明行政动作和机构立场，不自动证明所有地区、群体或对手按其意图行动。
\eventanchor{EQ-1648-LI-CHENGDONG-AFTERMATH}{「广东清将李成栋反清，广州一度…}顺治五年（1648年），「广东清将李成栋反清，广州一度转入永历朝控制」引发的善后、执行或地方回应构成可由不同材料追踪的后续节点。核心取证：《清世祖实录》顺治五年广东军报；《清史稿·李成栋传》；《广东通志》及永历朝诏敕辑录；相关地区的档案、志书、契约、报刊或对方材料。这一锚点追踪后续影响；不同地区、群体与时段的反应不能由战后或改革后的概括性记忆代替。来源保留：战后、改革后或条约后的记忆、地方记录和官方总结常具有事后选择性；后续影响须按地点、群体和时间分层，不作整体化推断。
\eventanchor{EQ-1649-NANCHANG-PRELUDE}{「清军围攻并攻取南昌，金声桓反…}顺治六年正月至三月（1649年），「清军围攻并攻取南昌，金声桓反叛失败」之前的条件、信息和争议已通过中央与地方材料显现，构成该事件可追踪的前史节点。核心取证：《清世祖实录》顺治六年正月至三月；《清史稿·金声桓传》；内阁大库顺治六年江西军需题本与《南昌府志》；同时期地方档案、地方志、日记或对方材料应按具体地区补检。这一锚点界定可回查的前史条件；条件、传闻或信息累积不应被写成已经证实的单一直接原因。来源保留：官修编年和地方材料的记录密度与立场并不相同；本行不将条件、传闻、呈报或后见解释直接等同为确定因果。
\eventanchor{EQ-1649-NANCHANG-DECISION}{围绕「清军围攻并攻取南昌，金声…}顺治六年正月至三月（1649年），围绕「清军围攻并攻取南昌，金声桓反叛失败」的命令、调度、照会或呈报形成独立的中央—地方行动文书链。核心取证：《清世祖实录》顺治六年正月至三月；《清史稿·金声桓传》；内阁大库顺治六年江西军需题本与《南昌府志》。这一锚点定位决策或调度动作；命令发出、地方接收、执行过程和实际结果仍须分别寻找证据。来源保留：奏折、上谕、部议、军机档或条约可证明行政动作和机构立场，不自动证明所有地区、群体或对手按其意图行动。
\eventanchor{EQ-1649-NANCHANG-AFTERMATH}{「清军围攻并攻取南昌，金声桓反…}顺治六年正月至三月（1649年），「清军围攻并攻取南昌，金声桓反叛失败」引发的善后、执行或地方回应构成可由不同材料追踪的后续节点。核心取证：《清世祖实录》顺治六年正月至三月；《清史稿·金声桓传》；内阁大库顺治六年江西军需题本与《南昌府志》；相关地区的档案、志书、契约、报刊或对方材料。这一锚点追踪后续影响；不同地区、群体与时段的反应不能由战后或改革后的概括性记忆代替。来源保留：战后、改革后或条约后的记忆、地方记录和官方总结常具有事后选择性；后续影响须按地点、群体和时间分层，不作整体化推断。
\subsection{1650—1659}
\eventanchor{EQ-1650-GUANGZHOU-PRELUDE}{「清军重新攻取广州，战后杀戮与…}顺治七年十一月（1650年），「清军重新攻取广州，战后杀戮与损毁规模成为广东地方记忆中的重大争议」之前的条件、信息和争议已通过中央与地方材料显现，构成该事件可追踪的前史节点。核心取证：《清世祖实录》顺治七年十一月；《清史稿·世祖本纪》；《广东通志》、广州地方笔记与内阁大库广东军报；同时期地方档案、地方志、日记或对方材料应按具体地区补检。这一锚点界定可回查的前史条件；条件、传闻或信息累积不应被写成已经证实的单一直接原因。来源保留：官修编年和地方材料的记录密度与立场并不相同；本行不将条件、传闻、呈报或后见解释直接等同为确定因果。
\eventanchor{EQ-1650-GUANGZHOU-DECISION}{围绕「清军重新攻取广州，战后杀…}顺治七年十一月（1650年），围绕「清军重新攻取广州，战后杀戮与损毁规模成为广东地方记忆中的重大争议」的命令、调度、照会或呈报形成独立的中央—地方行动文书链。核心取证：《清世祖实录》顺治七年十一月；《清史稿·世祖本纪》；《广东通志》、广州地方笔记与内阁大库广东军报。这一锚点定位决策或调度动作；命令发出、地方接收、执行过程和实际结果仍须分别寻找证据。来源保留：奏折、上谕、部议、军机档或条约可证明行政动作和机构立场，不自动证明所有地区、群体或对手按其意图行动。
\eventanchor{EQ-1650-GUANGZHOU-AFTERMATH}{「清军重新攻取广州，战后杀戮与…}顺治七年十一月（1650年），「清军重新攻取广州，战后杀戮与损毁规模成为广东地方记忆中的重大争议」引发的善后、执行或地方回应构成可由不同材料追踪的后续节点。核心取证：《清世祖实录》顺治七年十一月；《清史稿·世祖本纪》；《广东通志》、广州地方笔记与内阁大库广东军报；相关地区的档案、志书、契约、报刊或对方材料。这一锚点追踪后续影响；不同地区、群体与时段的反应不能由战后或改革后的概括性记忆代替。来源保留：战后、改革后或条约后的记忆、地方记录和官方总结常具有事后选择性；后续影响须按地点、群体和时间分层，不作整体化推断。
\eventanchor{EQ-1650-DORGON-DEATH-PRELUDE}{「摄政王多尔衮在古北口外行猎途…}顺治七年十二月（1650年），「摄政王多尔衮在古北口外行猎途中去世，摄政政治结构随之改组」之前的条件、信息和争议已通过中央与地方材料显现，构成该事件可追踪的前史节点。核心取证：《清世祖实录》顺治七年十二月；《清史稿·多尔衮传》；《清初内国史院满文档案》摄政王丧仪与遗命档；同时期地方档案、地方志、日记或对方材料应按具体地区补检。这一锚点界定可回查的前史条件；条件、传闻或信息累积不应被写成已经证实的单一直接原因。来源保留：官修编年和地方材料的记录密度与立场并不相同；本行不将条件、传闻、呈报或后见解释直接等同为确定因果。
\eventanchor{EQ-1650-DORGON-DEATH-DECISION}{围绕「摄政王多尔衮在古北口外行…}顺治七年十二月（1650年），围绕「摄政王多尔衮在古北口外行猎途中去世，摄政政治结构随之改组」的命令、调度、照会或呈报形成独立的中央—地方行动文书链。核心取证：《清世祖实录》顺治七年十二月；《清史稿·多尔衮传》；《清初内国史院满文档案》摄政王丧仪与遗命档。这一锚点定位决策或调度动作；命令发出、地方接收、执行过程和实际结果仍须分别寻找证据。来源保留：奏折、上谕、部议、军机档或条约可证明行政动作和机构立场，不自动证明所有地区、群体或对手按其意图行动。
\eventanchor{EQ-1650-DORGON-DEATH-AFTERMATH}{「摄政王多尔衮在古北口外行猎途…}顺治七年十二月（1650年），「摄政王多尔衮在古北口外行猎途中去世，摄政政治结构随之改组」引发的善后、执行或地方回应构成可由不同材料追踪的后续节点。核心取证：《清世祖实录》顺治七年十二月；《清史稿·多尔衮传》；《清初内国史院满文档案》摄政王丧仪与遗命档；相关地区的档案、志书、契约、报刊或对方材料。这一锚点追踪后续影响；不同地区、群体与时段的反应不能由战后或改革后的概括性记忆代替。来源保留：战后、改革后或条约后的记忆、地方记录和官方总结常具有事后选择性；后续影响须按地点、群体和时间分层，不作整体化推断。
\eventanchor{EQ-1651-SHUNZHI-PERSONAL-RULE-PRELUDE}{「顺治帝亲政，撤除摄政安排并重…}顺治八年（1651年），「顺治帝亲政，撤除摄政安排并重置皇帝周边的决策网络」之前的条件、信息和争议已通过中央与地方材料显现，构成该事件可追踪的前史节点。核心取证：《清世祖实录》顺治八年正月；《清史稿·世祖本纪》；《顺治朝内阁档》亲政诏令与任官题本；同时期地方档案、地方志、日记或对方材料应按具体地区补检。这一锚点界定可回查的前史条件；条件、传闻或信息累积不应被写成已经证实的单一直接原因。来源保留：官修编年和地方材料的记录密度与立场并不相同；本行不将条件、传闻、呈报或后见解释直接等同为确定因果。
\eventanchor{EQ-1651-SHUNZHI-PERSONAL-RULE-DECISION}{围绕「顺治帝亲政，撤除摄政安排…}顺治八年（1651年），围绕「顺治帝亲政，撤除摄政安排并重置皇帝周边的决策网络」的命令、调度、照会或呈报形成独立的中央—地方行动文书链。核心取证：《清世祖实录》顺治八年正月；《清史稿·世祖本纪》；《顺治朝内阁档》亲政诏令与任官题本。这一锚点定位决策或调度动作；命令发出、地方接收、执行过程和实际结果仍须分别寻找证据。来源保留：奏折、上谕、部议、军机档或条约可证明行政动作和机构立场，不自动证明所有地区、群体或对手按其意图行动。
\eventanchor{EQ-1651-SHUNZHI-PERSONAL-RULE-AFTERMATH}{「顺治帝亲政，撤除摄政安排并重…}顺治八年（1651年），「顺治帝亲政，撤除摄政安排并重置皇帝周边的决策网络」引发的善后、执行或地方回应构成可由不同材料追踪的后续节点。核心取证：《清世祖实录》顺治八年正月；《清史稿·世祖本纪》；《顺治朝内阁档》亲政诏令与任官题本；相关地区的档案、志书、契约、报刊或对方材料。这一锚点追踪后续影响；不同地区、群体与时段的反应不能由战后或改革后的概括性记忆代替。来源保留：战后、改革后或条约后的记忆、地方记录和官方总结常具有事后选择性；后续影响须按地点、群体和时间分层，不作整体化推断。
\eventanchor{EQ-1652-DINGGUO-GUILIN-PRELUDE}{「李定国攻克桂林，孔有德自尽，…}永历六年／顺治九年（1652年），「李定国攻克桂林，孔有德自尽，清在广西的军事据点受重创」之前的条件、信息和争议已通过中央与地方材料显现，构成该事件可追踪的前史节点。核心取证：《清世祖实录》顺治九年广西军报；《清史稿·孔有德传》；《永历实录》与《广西通志》；同时期地方档案、地方志、日记或对方材料应按具体地区补检。这一锚点界定可回查的前史条件；条件、传闻或信息累积不应被写成已经证实的单一直接原因。来源保留：官修编年和地方材料的记录密度与立场并不相同；本行不将条件、传闻、呈报或后见解释直接等同为确定因果。
\eventanchor{EQ-1652-DINGGUO-GUILIN-DECISION}{围绕「李定国攻克桂林，孔有德自…}永历六年／顺治九年（1652年），围绕「李定国攻克桂林，孔有德自尽，清在广西的军事据点受重创」的命令、调度、照会或呈报形成独立的中央—地方行动文书链。核心取证：《清世祖实录》顺治九年广西军报；《清史稿·孔有德传》；《永历实录》与《广西通志》。这一锚点定位决策或调度动作；命令发出、地方接收、执行过程和实际结果仍须分别寻找证据。来源保留：奏折、上谕、部议、军机档或条约可证明行政动作和机构立场，不自动证明所有地区、群体或对手按其意图行动。
\eventanchor{EQ-1652-DINGGUO-GUILIN-AFTERMATH}{「李定国攻克桂林，孔有德自尽，…}永历六年／顺治九年（1652年），「李定国攻克桂林，孔有德自尽，清在广西的军事据点受重创」引发的善后、执行或地方回应构成可由不同材料追踪的后续节点。核心取证：《清世祖实录》顺治九年广西军报；《清史稿·孔有德传》；《永历实录》与《广西通志》；相关地区的档案、志书、契约、报刊或对方材料。这一锚点追踪后续影响；不同地区、群体与时段的反应不能由战后或改革后的概括性记忆代替。来源保留：战后、改革后或条约后的记忆、地方记录和官方总结常具有事后选择性；后续影响须按地点、群体和时间分层，不作整体化推断。
\eventanchor{EQ-1652-DINGGUO-HENGZHOU-PRELUDE}{「李定国在衡州击败尼堪军，清军…}永历六年／顺治九年夏（1652年），「李定国在衡州击败尼堪军，清军西南攻势暂受挫」之前的条件、信息和争议已通过中央与地方材料显现，构成该事件可追踪的前史节点。核心取证：《清世祖实录》顺治九年湖南军报；《清史稿·尼堪传》；《永历实录》与《衡州府志》；同时期地方档案、地方志、日记或对方材料应按具体地区补检。这一锚点界定可回查的前史条件；条件、传闻或信息累积不应被写成已经证实的单一直接原因。来源保留：官修编年和地方材料的记录密度与立场并不相同；本行不将条件、传闻、呈报或后见解释直接等同为确定因果。
\eventanchor{EQ-1652-DINGGUO-HENGZHOU-DECISION}{围绕「李定国在衡州击败尼堪军，…}永历六年／顺治九年夏（1652年），围绕「李定国在衡州击败尼堪军，清军西南攻势暂受挫」的命令、调度、照会或呈报形成独立的中央—地方行动文书链。核心取证：《清世祖实录》顺治九年湖南军报；《清史稿·尼堪传》；《永历实录》与《衡州府志》。这一锚点定位决策或调度动作；命令发出、地方接收、执行过程和实际结果仍须分别寻找证据。来源保留：奏折、上谕、部议、军机档或条约可证明行政动作和机构立场，不自动证明所有地区、群体或对手按其意图行动。
\eventanchor{EQ-1652-DINGGUO-HENGZHOU-AFTERMATH}{「李定国在衡州击败尼堪军，清军…}永历六年／顺治九年夏（1652年），「李定国在衡州击败尼堪军，清军西南攻势暂受挫」引发的善后、执行或地方回应构成可由不同材料追踪的后续节点。核心取证：《清世祖实录》顺治九年湖南军报；《清史稿·尼堪传》；《永历实录》与《衡州府志》；相关地区的档案、志书、契约、报刊或对方材料。这一锚点追踪后续影响；不同地区、群体与时段的反应不能由战后或改革后的概括性记忆代替。来源保留：战后、改革后或条约后的记忆、地方记录和官方总结常具有事后选择性；后续影响须按地点、群体和时间分层，不作整体化推断。
\eventanchor{EQ-1653-ZHENG-XIAMEN-PRELUDE}{「郑成功以厦门、金门为基地整合…}顺治十年（1653年），「郑成功以厦门、金门为基地整合海上军政与贸易网络，持续攻击福建沿海清军据点」之前的条件、信息和争议已通过中央与地方材料显现，构成该事件可追踪的前史节点。核心取证：《清世祖实录》顺治十年福建海防军报；《清史稿·郑成功传》；《厦门志》与郑氏文书辑录；同时期地方档案、地方志、日记或对方材料应按具体地区补检。这一锚点界定可回查的前史条件；条件、传闻或信息累积不应被写成已经证实的单一直接原因。来源保留：官修编年和地方材料的记录密度与立场并不相同；本行不将条件、传闻、呈报或后见解释直接等同为确定因果。
\eventanchor{EQ-1653-ZHENG-XIAMEN-DECISION}{围绕「郑成功以厦门、金门为基地…}顺治十年（1653年），围绕「郑成功以厦门、金门为基地整合海上军政与贸易网络，持续攻击福建沿海清军据点」的命令、调度、照会或呈报形成独立的中央—地方行动文书链。核心取证：《清世祖实录》顺治十年福建海防军报；《清史稿·郑成功传》；《厦门志》与郑氏文书辑录。这一锚点定位决策或调度动作；命令发出、地方接收、执行过程和实际结果仍须分别寻找证据。来源保留：奏折、上谕、部议、军机档或条约可证明行政动作和机构立场，不自动证明所有地区、群体或对手按其意图行动。
